\chapter{Lista 9.}

\noindent\textbf{Przypomnienie:}

Niech $X = \{x_{i}: i \in I\} $ będzie dowolnym zbiorem. Przez $\mathbb{L}_{X}$ oznaczamy zbiór wszystkich słów nad alfabetem $X$, a przez  $N_{X}$ zbiór słów nieskracalnych. Przypomnijmy, że grupa $\mathbb{F}_{X}$ była zdefiniowana jako zbiór klas abstrakcji relacji $\sim  $ na $\mathbb{L}_{X}$ z działaniem konkatenacji na reprezentantach. Grupa wolna o wolnym zbiorze generatorów była zdefiniowana jako dowolna grupa $G \supseteq X$, dla której istnieje izomorfizm $\varphi : G \to \mathbb{F}_{X} $, taki że $\varphi \left( x \right) = \left[ x \right]  $ dla każdego $x \in X$. Utożsamiając  $x$ z $\left[ x \right] $, $\mathbb{F}_{X}$ staje się grupą wolną o wolnym zbiorze generatorów $X$. Dla wygody czasami dowolną grupę wolną o wolnym zbiorze generatorów $X$ oznacza się przez $\mathbb{F}_{X}$ (chociaż jest ona zdefiniowana z dokładnością do izomorfizmu nad X); z kontekstu powinno być jasne, czy $\mathbb{F}_{X}$ oznacza dowolną grupę wolną, czy jej konkretną realizację za pomocą pierwotnej definicji $\mathbb{F}_{X}$ jako zbioru klas abstrakcji relacji $\sim$.
\begin{problem}[1]
	Niech $\rho$ będzie funkcją ze zbioru $\mathbb{L}_{X}$ w zbiór $N_{X}$, polegającą na kolejnym usuwaniu znajdującego się najbardziej na prawo podsłowa postaci $x_{i}^{\varepsilon}x_{i}^{-\varepsilon}$, gdzie $i \in I$ oraz $\varepsilon \in \{-1,1\} $. Udowodnić następujące własności funkcji $\rho$, gdzie $u,v \in \mathbb{L}_{X}$ oraz $\varepsilon \in \{-1,1\}$.
	\begin{enumerate}[label=(\alph*)]
		\item $\rho\left( u \right) \sim u $.
		\item $\rho\left( u \right) = u \leftrightarrow u \text{ jest nieskracalne}$.
		\item $\rho\left( uv \right) = \rho\left( u\rho\left( v \right)  \right) $.
		\item $\rho\left( x_{i}^{\varepsilon}x_{i}^{-\varepsilon }u \right) = \rho\left( u \right) $.
		\item $\rho\left( ux_{i}^{\varepsilon }x_{i}^{-\varepsilon } v\right) \rho\left( uv \right) $.
		\item $\rho\left( uv \right) = \rho\left( \rho\left( u \right) \rho\left( v \right)\right)$. (Wskazówka: indukcja względem $\left| u \right|.$
	\end{enumerate}
\end{problem} 
\begin{solution}

\end{solution} 
\begin{problem}[2]
	
\end{problem} 
